\documentclass{article}
\usepackage{tcee-repaso}
\tcetema[La coordinación internacional de políticas económicas]{3.B.19}{La coordinación internacional de políticas económicas. Aspectos teóricos y prácticos. El G-20, la OCDE y otros foros internacionales.}
\tcefecha{13/11/2026}
\begin{document}
\cabecerarepaso
\begin{columnas}

% ==========================================================================
\bloque{Esquema del cante}
\esquemacante{%
\apartadocante{Introducción: «cada uno en su casa y Dios en la de todos»}{2,5}
\apartadocante{I. Análisis teórico}{12,5}
\subcante{1. Condición necesaria: desbordamientos}{8,5}
\subcante{\hspace*{0.8em}Canales}{0,5}
\subcante{\hspace*{0.8em}\autor{Mundell}-\autor{Fleming} con 2 países}{3,5}
\subcante{\hspace*{0.8em}\autor{Corsetti} y \autor{Pesenti}}{4,5}
\subcante{2. Condición suficiente: teoría de juegos}{4}
\subcante{\hspace*{0.8em}\autor{Hamada}}{2,5}
\subcante{\hspace*{0.8em}La repetición}{1}
\subcante{\hspace*{0.8em}Evidencia y valoración}{0,5}
\apartadocante{II. La cooperación en la historia reciente}{6}
\subcante{1. Ejemplos de cooperación}{4}
\subcante{2. Obstáculos}{1,5}
\subcante{3. Formas}{0,5}
\apartadocante{III. Economía política: mecanismos}{8}
\subcante{1. Por qué hacen falta foros}{1,5}
\subcante{2. G-20, OCDE y otros foros}{6,5}
\apartadocante{Conclusión}{1}}
Funciones clave: las 3 curvas de \autor{Corsetti}-\autor{Pesenti}; la pérdida de \autor{Hamada}; la condición del gatillo. Gráficos: fiscal con 2 países, \autor{Corsetti}-\autor{Pesenti}, \autor{Hamada}; figura: escalera de la cooperación; tabla: línea del tiempo.

\emph{Pizarra}: introducción y conclusión sin pizarra. Siempre los mismos 2 países: A (zona del euro) y B (EE.\,UU.). I: 5 canales en una columna estrecha; la fiscal con 2 paneles por partes; recuadro con las 3 curvas y \autor{Corsetti}-\autor{Pesenti} por capas ($O\to R\to F$; $H$ y $L$); \autor{Hamada} por capas; la matriz de la locomotora y el gatillo, en una esquina hasta el final. II: flecha del tiempo (1944 … 2025) que crece con los episodios; sobre la matriz se tachan los rótulos (Expandir → Consolidar; Gastar → Contener/Depreciar). III: la escalera, peldaño a peldaño, con año y miembros de cada foro.

% ==========================================================================
\bloque[2,5]{Introducción}
\begin{itemize}
\item \textbf{Enganche}: «cada uno en su casa y Dios en la de todos». Johannesburgo, 22-23/11/2025: 1.ª cumbre del G-20 en África, \emph{sin} EE.\,UU. (que preside en 2026). Londres, 2009: la mayor expansión fiscal concertada. ¿Basta con poner orden en casa o hay problemas que solo se arreglan entre todos?
\item \textbf{Relevancia}: la Fed mueve crédito y activos en todo el mundo (\autor{Rey}) y, por el dólar, los precios del comercio entre terceros (\autor{Gopinath}); el FMI pide reequilibrar a la vez (ESR 2026); España: invitado permanente del G-20, fundador de la OCDE, miembro del FSB y del euro \vertema{3.B.45}; une macro abierta, juegos y economía política.
\item \textbf{Contextualización}: \concepto{cooperación} = escalera (información → acuerdo puntual → reglas); \concepto{coordinación} = cada país cambia su política, de mutuo acuerdo, por sus efectos sobre los demás (de \autor{Nash} a solución negociada). Puntual o permanente; demanda u oferta. 3 etapas de las ideas: \autor{Friedman} (1953) en contra; 70-80, modelos keynesianos y juegos a favor (Bonn, Plaza, Louvre); NOEM escéptica (\autor{Obstfeld}-\autor{Rogoff}, 2002), y 2008 la devuelve con el G-20. $E$ directo: $E\uparrow$ = depreciación.
\item \textbf{Problemática}: ¿por qué y por qué canales? ¿Cuándo es deseable y por qué se ve tan poco? ¿Reglas o acuerdos, en qué foros?
\item \textbf{Estructura}: \hyperref[rep:I]{I}, teoría (necesaria y suficiente); \hyperref[rep:II]{II}, historia; \hyperref[rep:III]{III}, foros.
\end{itemize}

% ==========================================================================
\bloque[12,5]{I. Análisis teórico\label{rep:I}}
\begin{ideaclave}
Condición necesaria: hay desbordamientos, pero su signo depende del modelo y de cómo se fijan los precios. Condición suficiente: coordinar compensa; sin coordinar el resultado es ineficiente, pero la solución cooperativa no es un equilibrio. Desbordamientos sin ganancias de coordinar son posibles (LCP).
\end{ideaclave}

\sub{I.1. Condición necesaria: los desbordamientos\label{rep:I.1}}
\begin{itemize}
\item 2 motivos (\autor{Bergin}, 2008): \concepto{efectos desbordamiento} (\emph{spillovers}) = externalidades \vertema{3.A.23}; reparto internacional del riesgo.
\item Fiscal: cada uno espera a que tire otro (bien público); monetaria: \concepto{depreciaciones competitivas}. Peso según la integración.
\item 5 canales: \emph{renta} (locomotora), \emph{tipo de interés}, \emph{tipo de cambio}, \emph{precios relativos} (relación real de intercambio, solo con microfundamentos), \emph{financiero}.
\item \concepto{Ciclo financiero global}: la Fed mueve precios de activos, apalancamiento y crédito en todo el mundo; un \emph{dilema}, no un trilema (\autor{Rey}, 2013; \autor{Miranda-Agrippino}-\autor{Rey}, 2020) \vertema{3.B.15}; el BCE matiza (\autor{Habib}-\autor{Venditti}, 2018).
\item Las políticas de oferta tienen desbordamientos más claros y más institucionalizados (OMC, OCDE, clima).
\end{itemize}

\textbf{\autor{Mundell}-\autor{Fleming} con 2 países} (\autor{Mundell}, 1964). \emph{Supuestos}: $P=P^*=1$, expectativas estáticas, sustitutos perfectos, $\beta_K\to\infty$, \autor{Marshall}-\autor{Lerner}; 2 países grandes con tipo flexible; $i=i^*$ y renta exterior endógenos.
\pareja{mundell-fleming-dos-paises-fiscal}{Fiscal en A: sube $i^*$, B se deprecia y su renta sube (locomotora)}{mundell-fleming-dos-paises}{Monetaria en A: B se aprecia y su renta baja (empobrecimiento)}
\begin{itemize}
\item \emph{Monetaria en A}: LM$\to$, $i\downarrow$, A se deprecia; en B dominan la apreciación: renta $\downarrow$: \concepto{empobrecimiento del vecino} (\emph{beggar-thy-neighbour}).
\item \emph{Fiscal en A}: IS$\to$, $i\uparrow$, A se aprecia; en B las 2 fuerzas suman: renta $\uparrow$: \concepto{efecto locomotora}.
\item Con tipo fijo: monetaria sube las 2 rentas; fiscal, ambigua en B. A pequeño: gana eficacia la monetaria; A grande: la fiscal.
\item En 4 ecuaciones: sumando IS y LM sale el IS-LM del mundo ($Y_A+Y_B$, $i$); con $i$ fijado, $Y_B=(M_B+h\,i)/k$: monetaria de A ($i\downarrow$) $\Rightarrow Y_B\downarrow$; fiscal ($i\uparrow$) $\Rightarrow Y_B\uparrow$; con LM iguales, $\Delta Y_A=\Delta Y_B=(h/k)\Delta i$.
\item Evidencia: la expansión de la Fed \emph{sube} la producción fuera (\autor{Kim}, 2001; \autor{Dedola} et al., 2017). Valoración: estático, sin bienestar, \autor{Lucas} (1976); políticas exógenas.
\end{itemize}

\textbf{\autor{Corsetti} y \autor{Pesenti} (2009)}. NOEM desde el \emph{Redux} (\autor{Obstfeld}-\autor{Rogoff}, 1995): optimización, competencia monopolística, rigideces; aporta la relación de intercambio y un criterio de bienestar. $U=\ln C-\chi\ell$; $\sigma>1$; precios fijados un periodo; mercados completos; banco central benevolente con tono $\mu$.
\funcionclave{\text{AD: }C=\frac{\mu}{P},\quad \text{AS: }C=Z\ell\tau,\quad \text{NR: }\overline{\ell}=\frac{\sigma-1}{\sigma\chi},\quad E=\frac{\mu}{\mu^*}}
\graficorepaso[0.95\linewidth]{corsetti-pesenti}{a) $Z\uparrow$ en A: $O\to R$ (rígidos), $F$ con la respuesta monetaria; b) B en $H$ (PCP) o en $L$ (LCP)}
\begin{itemize}
\item \concepto{Traspaso} (\emph{pass-through}): \emph{PCP} completo; \emph{LCP} nulo a corto; \emph{dollar pricing} asimétrico.
\item PCP: $\tau\downarrow$ en A, $\tau_B\uparrow$; B a $H$: desbordamiento \emph{positivo} (al revés que M-F; canal = relación de intercambio). Sin sobrerreacción.
\item LCP: cambian los márgenes; $\tau_B\downarrow$; B a $L$, más trabajo: \emph{empobrece} al vecino; sin desvío del gasto; sobrerreacción con \emph{pricing to market} (\autor{Betts}-\autor{Devereux}).
\item ¿Ganancias de coordinar? PCP: ninguna; LCP: \emph{tampoco} (los 2 eligen la misma política; \autor{Devereux}-\autor{Engel}, 2003); dólar: sí, pero unilaterales del emisor (\autor{Gopinath}, moneda dominante).
\item Extensiones: reforma estructural ($\sigma\uparrow$) = desbordamiento positivo de oferta (gráfico abajo); reparto del riesgo (1 instrumento, 2 distorsiones); bienes intermedios; no comerciables. Evidencia: ganancias pequeñas, hasta 2,2\,\% del consumo (\autor{Tchakarov}, 2004).
\end{itemize}

\graficorepaso[0.95\linewidth]{corsetti-pesenti-reforma}{Reforma en A ($\sigma\uparrow$): A a $O'$, B a $O_B'$: los 2 ganan}
\graficorepaso[0.7\linewidth]{corsetti-pesenti-dolar}{Expansión del emisor del dólar vista desde A: a $D$ (consume como en $H$, trabaja más)}
\begin{itemize}
\item Con dólar, el emisor (B) queda aislado: sus precios de importación están fijos en su moneda; A recibe todo el efecto.
\end{itemize}

\sub{I.2. Condición suficiente: la teoría de juegos\label{rep:I.2}}
\textbf{El diagrama de \autor{Hamada}} (1974, 1976; versión de \autor{Corden}, 1994). Quién le pone el cascabel al gato: los 2 saben lo que conviene y cada uno espera al otro. 2 países, instrumento $M_i$; disyuntiva de \autor{Phillips}; $E=f(M_A/M_B)$; objetivo común $E_0$; $\gamma_B>\gamma_A$; juego estático.
\funcionclave{L_i=(U_i-U^*_i)^2+\gamma_i(\pi_i-\pi^*_i)^2,\qquad i\in\{A,B\}}
\graficorepaso[0.8\linewidth]{hamada}{Reacciones (vertical por $F_A$, horizontal por $F_B$), Nash $N$, lente, curva de contrato $F_AF_B$ y núcleo sobre $E_0$}
\begin{itemize}
\item Rayos = tipos de cambio; círculos alrededor del \concepto{punto de felicidad}. $N$ no es Pareto: lente; en $N$, A demasiado depreciada.
\item Cooperativa: $\min\omega L_A+(1-\omega)L_B$; \concepto{curva de contrato}; \concepto{núcleo} = la parte dentro de la lente, sobre $E_0$.
\item Pero no es un equilibrio: incentivo a ser \concepto{polizón} (\emph{free rider}); dilema del prisionero continuo. \autor{Stackelberg}: con círculos, no cambia nada.
\item Evidencia: \autor{Oudiz}-\autor{Sachs} (1984): medio punto del PNB al año para EE.\,UU.; poco comercio (1,6\,\% del PNB en 1982). «Monstruo del lago Ness» (\autor{Blanchard} et al., 2013).
\end{itemize}

\textbf{La repetición}. Locomotora (\autor{Frankel}, 2016): (5,5), (0,6), (6,0), (1,1); no expandir, dominante. Juego repetido sin final, \concepto{estrategia del gatillo} (\autor{James W. Friedman}, 1971):
\funcionclave{\beta\geq\frac{u^{\mathrm D}-u^{\mathrm C}}{u^{\mathrm D}-u^{\mathrm N}}\quad\Rightarrow\quad \beta\geq\tfrac{6-5}{6-1}=\tfrac15}
\begin{itemize}
\item \emph{Folk theorem} \vertema{3.A.14}; $\beta$ = paciencia de los gobiernos (elecciones, reputación). No basta: pocas cumbres, ruido, castigo poco creíble $\Rightarrow$ reglas e instituciones.
\end{itemize}

% ==========================================================================
\bloque[6]{II. La cooperación en la historia reciente\label{rep:II}}
\begin{ideaclave}
Se ha coordinado poco y casi siempre en crisis; cada episodio es uno de los juegos, y el desacuerdo suele ser sobre el signo del desbordamiento, no sobre cooperar. Lo que dura son las reglas.
\end{ideaclave}
\sub{II.1. Ejemplos de cooperación\label{rep:II.1}}
\begin{itemize}
\item \emph{Fiscal}: Bonn 1978 (Alemania hasta 1\,\% del PNB, Francia 0,5, Japón +1,5 puntos; Alemania se arrepintió); Londres, 2/4/2009 (\dato{5 billones \$ + 1,1; FMI ×3 hasta 750.000 M}{G-20}) \vertema{3.A.40}; Toronto 2010 (déficits a la mitad antes de 2013) $\Rightarrow$ \emph{disciplina fiscal}: PEC reforzado (2011) y reformado (2024) \vertema{3.B.45}; Alemania veía disciplina y EE.\,UU. locomotora.
\item \emph{Monetaria}: QE de la Fed (11/2008, 11/2010, 9/2012), BoJ (4/2013), BCE (3/2015); «guerra de divisas» (\autor{Mantega}, 2010), «tsunami monetario» (\autor{Rousseff}, 2012); G-7, 12/2/2013: objetivos internos con instrumentos internos. Matiz: en el límite inferior, locomotora (\autor{Eichengreen}-\autor{Sachs}, 1985). Bajada conjunta 8/10/2008; \emph{swaps} de la Fed (2008, 2020).
\item \emph{Tipos de cambio}: Plaza (22/9/1985, G-5), Tokio (5/1986, indicadores), Louvre (22/2/1987, G-6) = \concepto{batalla de los sexos}.
\item \emph{¿Quién ajusta?} Plan de \autor{Keynes} (1943): cobrar también al acreedor; ganó \autor{White} \vertema{3.B.20}. Hoy: ESR del FMI, evaluación mutua del G-20, procedimiento de desequilibrios.
\end{itemize}
\sub{II.2. Obstáculos\label{rep:II.2}}
\begin{itemize}
\item Económicos: polizón; contraproducente (\autor{Rogoff}, 1985: coordinar solo a bancos centrales sin credibilidad sube la inflación); incertidumbre sobre el modelo (\autor{Frankel}-\autor{Rockett}, 1988); preferencias ($\gamma_i$); reparto.
\item Políticos: soberanía; juegos de 2 niveles (\autor{Putnam}, 1988); mandatos nacionales. A favor: la moneda muy depreciada alimenta el proteccionismo (\autor{Broz}-\autor{Frieden}, 2013).
\end{itemize}
\sub{II.3. Formas\label{rep:II.3}}
\begin{itemize}
\item \concepto{Reglas} (patrón oro, Bretton Woods, SME, reglas fiscales: credibilidad, poca flexibilidad) frente a \concepto{acuerdos} (Plaza, Londres: flexibles, duran poco). Con líder hegemónico o entre iguales.
\end{itemize}

% ==========================================================================
\bloque[8]{III. Economía política: mecanismos\label{rep:III}}
\begin{ideaclave}
Un foro cambia el juego: lo repite, da información, hace visibles los incumplimientos y liga asuntos. G-20, foro de las crisis y la demanda; OCDE, de las normas y la oferta.
\end{ideaclave}
\sub{III.1. Por qué hacen falta foros\label{rep:III.1}}
\begin{itemize}
\item Diseño: \emph{amplitud} (compensaciones entre asuntos), \emph{profundidad} (vigilar mejor), \emph{tamaño} (G-7 barato; G-20 internaliza más). \emph{Clubes}: ventaja excluible (\autor{Nordhaus}, 2015) \vertema{3.B.31}.
\end{itemize}
\graficorepaso[0.8\linewidth]{escalera-cooperacion}{Figura: información (OCDE, FMI) → coordinación puntual (G-7, G-20) → reglas (OMC, UE)}
\sub{III.2. G-20, OCDE y otros foros\label{rep:III.2}}
\begin{itemize}
\item \emph{G-20}: 1999 (ministros, Berlín); líderes desde Washington (11/2008); «foro principal» (Pittsburgh, 2009). 19 países + UE + Unión Africana (2023); \dato{$\approx$85\,\% PIB, 75\,\% comercio, 2/3 población}{G-20}; España invitado permanente. Informal, consenso, \emph{troika}, \emph{sherpas} y vía financiera. Hizo: 2009, FSB, Marco para un Crecimiento Fuerte, Sostenible y Equilibrado con evaluación mutua, Basilea~III, 2 pilares, ISSD (\dato{>12.000 M\$, casi 50 países}{G-20; BM}) y Marco Común \vertema{3.B.29}. Límites: no vinculante, pierde fuelle fuera de crisis, dividido (Johannesburgo 2025). Presidencias: Sudáfrica 2025, EE.\,UU. 2026 (Miami, 14-15/12), Reino Unido 2027.
\item \emph{OCDE}: sucede a la OECE (1948, Plan Marshall); Convenio de París 14/12/1960, 20 firmantes, España fundadora; \dato{38 miembros, 8 candidatos}{OCDE}; Consejo, comités, Secretaría (\autor{Cormann}). \concepto{Presión entre iguales}; \emph{Economic Outlook}, estudios, \emph{Going for Growth}; decisiones (Código de capitales), recomendaciones, Consenso de créditos a la exportación \vertema{4.A.24}, CAD; BEPS y 2 pilares (mínimo 15\,\%, 750 M€) \vertema{4.B.18}.
\item \emph{Otros}: G-7 (Rusia suspendida desde 2014; Club del Clima, 2022); FSB (4/2009, asociación suiza desde 2013, \autor{Bailey}); FMI (art.~IV, WEO, ESR, alerta temprana con el FSB, visión institucional 2012/2022); BPI y Comité de Basilea; OMC; Club de París; UE (la más intensa); RIC, Davos, ASEM.
\end{itemize}

% ==========================================================================
\bloque[1]{Conclusión}
\begin{itemize}
\item \textbf{Recapitulación}: necesaria si hay desbordamientos; deseable porque \autor{Nash} es ineficiente (\autor{Hamada}); difícil porque la cooperación no es un equilibrio; se coordina en crisis y con reglas.
\item \textbf{Relevancia}: desequilibrios crecientes, foros divididos, ciclo financiero global.
\item \textbf{Extensiones}: \vertema{3.B.15} \vertema{3.B.16} \vertema{3.B.45} \vertema{3.B.31} \vertema{4.B.18} \vertema{3.B.29} \vertema{3.B.22}.
\item \textbf{Opinión}: el «monstruo del lago Ness»: se ve poco, pero aparece cuando más importa, y sus reglas duran.
\item \textbf{Idea final}: cada uno en su casa…, pero cuando la casa de uno ocupa media calle, la de todos también es suya.
\end{itemize}

% ==========================================================================
\bloque{Definiciones}
\begin{itemize}
\item \concepto{Coordinación}: cambiar la política propia de mutuo acuerdo para internalizar sus efectos sobre otros. \concepto{Cooperación}: cualquier peldaño de la escalera.
\item \concepto{Relación real de intercambio} ($\tau$): consumo que compra una unidad de lo producido.
\item \concepto{PCP, LCP, dollar pricing}: precio en la moneda del productor, del comprador o de un país (dólar).
\item \concepto{Punto de felicidad}: pérdida nula. \concepto{Núcleo}: acuerdos que mejoran a los 2 frente a $N$.
\item \concepto{Estrategia del gatillo}: cooperar hasta que el otro se desvíe; luego, \autor{Nash} para siempre.
\item \concepto{Presión entre iguales}: evaluación de cada país por los demás (OCDE).
\item \concepto{Ciclo financiero global}: comovimiento mundial de precios de activos, flujos brutos y crédito, movido por la Fed.
\item \concepto{Paradigma de la moneda dominante}: el comercio se factura en dólares; el emisor queda aislado (\autor{Gopinath} et al., 2020).
\item \concepto{Juego de 2 niveles}: cada gobierno negocia a la vez fuera y dentro (\autor{Putnam}, 1988).
\item \concepto{Club}: acuerdo con ventaja excluible para sostener la participación (\autor{Nordhaus}, 2015).
\end{itemize}

% ==========================================================================
\bloque{Datos clave}
\begin{itemize}
\item Londres 2009: \dato{5 billones \$; FMI hasta 750.000 M\$}{G-20}. Toronto: \dato{déficits a la mitad antes de 2013}{G-20}.
\item G-20: \dato{$\approx$85\,\% PIB, 75\,\% comercio, 2/3 población}{G-20}. OCDE: \dato{38 miembros, 8 candidatos}{OCDE, 2026}.
\item Pilar 2: \dato{mínimo del 15\,\% a grupos con $\geq$750 M€}{G-20, 2021}.
\item Superávit 2025 (\autor{Keynes}: el ajuste simétrico sigue pendiente): \dato{saldo público Alemania $-2{,}9$, Países Bajos $-1{,}6$, Corea $-2{,}3$\,\% del PIB}{FMI, WEO 10/2026}; freno de la deuda alemán reformado (3/2025, fondo de 500.000 M€).
\end{itemize}
\graficorepaso[0.9\linewidth]{saldos-corrientes-fiscales-superavit}{Alemania, Países Bajos y Corea, 2019 y 2025: siguen con superávit exterior, ya con déficit público}

% ==========================================================================
\bloque{\autor{Mundell}-\autor{Fleming} con 2 países: resumen}
{\small\noindent\begin{tabular}{@{}l|l|l|l@{}}
\textbf{Régimen} & \textbf{Política} & \textbf{Renta de A} & \textbf{Renta de B}\\ \hline
Flexible & monetaria & $\uparrow$ (más si A pequeño) & $\downarrow$ (empobrece)\\
Flexible & fiscal & $\uparrow$ (más si A grande) & $\uparrow$ (locomotora)\\
Fijo & monetaria & $\uparrow$ & $\uparrow$\\
Fijo & fiscal & $\uparrow$ & ambiguo\\
\end{tabular}\par}

% ==========================================================================
\bloque{Línea del tiempo de la coordinación}
{\small\noindent\begin{tabular}{@{}l|p{0.74\linewidth}@{}}
\textbf{Año} & \textbf{Hito y lectura}\\ \hline
1944 & Bretton Woods; ajuste solo del deudor\\
1978 & Bonn (G-7): locomotora de Alemania y Japón\\
1985 & Plaza (G-5): depreciar el dólar\\
1986-87 & Tokio y Louvre (G-6): indicadores y estabilizar\\
1999 & G-20 de ministros y FEF: voz a los emergentes\\
2008-09 & Bajada conjunta; Washington, Londres, Pittsburgh: locomotora mundial; FSB\\
2010 & Toronto: disciplina fiscal\\
2013 & G-7: alto el fuego en la guerra de divisas\\
2020-21 & \emph{Swaps}, ISSD, Marco Común, 2 pilares\\
2025 & Johannesburgo sin EE.\,UU.: foros divididos\\
\end{tabular}\par}

% ==========================================================================
\bloque{Los 3 juegos de \autor{Frankel} (mismos pagos)}
{\small\noindent\begin{tabular}{@{}l|l|l@{}}
\textbf{Juego} & \textbf{Cooperar} & \textbf{Dominante}\\ \hline
Locomotora & expandir & no expandir\\
Disciplina fiscal & consolidar & gastar\\
Guerra de divisas & contener & depreciar\\
Plaza-Louvre & \multicolumn{2}{l}{batalla de los sexos: 2 equilibrios}\\
\end{tabular}\par}

% ==========================================================================
\bloque{Los foros de un vistazo}
{\small\noindent\begin{tabular}{@{}>{\raggedright\arraybackslash}p{0.14\linewidth}|>{\raggedright\arraybackslash}p{0.19\linewidth}|>{\raggedright\arraybackslash}p{0.23\linewidth}|>{\raggedright\arraybackslash}p{0.26\linewidth}@{}}
\textbf{Foro} & \textbf{Origen} & \textbf{Miembros} & \textbf{Coordina}\\ \hline
G-20 & 1999; líderes 2008 & 19 + UE + UA & crisis, demanda, finanzas\\
OCDE & 1960 (OECE, 1948) & 38; 8 candidatos & normas, oferta, fiscalidad\\
G-7 & foro de los 70-80 & 7 + UE (Rusia, suspendida) & avanzadas; Club del Clima\\
FSB & 4/2009 (FEF, 1999) & G-20 y otros, España & reguladores\\
FMI & Bretton Woods (1944) & con tratado y recursos & art.~IV, WEO, ESR\\
OMC & — & — & comercio, con reglas y solución de diferencias\\
Club de París & — & acreedores oficiales & deuda bilateral, con el Marco Común\\
UE & — & Estados miembros & reglas fiscales, Semestre, euro\\
\end{tabular}\par}

% ==========================================================================
\bloque{Autores y fechas\label{rep:autores}}
\hito{1943}{Plan de \autor{Keynes}}
\hito{1953}{\autor{Friedman}: flexibles aíslan}
\hito{1964}{\autor{Mundell}: 2 países}
\hito{1971}{\autor{J. W. Friedman}: gatillo}
\hito{1976}{\autor{Hamada} (1974)}
\hito{1978}{Cumbre de Bonn}
\hito{1984}{\autor{Oudiz}-\autor{Sachs}}
\hito{1985}{Plaza; \autor{Rogoff}: contraproducente}
\hito{1987}{Louvre}
\hito{1988}{\autor{Frankel}-\autor{Rockett}; \autor{Putnam}}
\hito{1995}{\emph{Redux}}
\hito{1999}{G-20 de ministros; FEF}
\hito{2003}{\autor{Devereux}-\autor{Engel}}
\hito{2009}{Londres; FSB; \autor{Corsetti}-\autor{Pesenti}}
\hito{2013}{G-7 sobre divisas; \autor{Rey}}
\hito{2021}{2 pilares (Roma)}
\hito{2025}{Johannesburgo sin EE.\,UU.}

% ==========================================================================
\bloque{Preguntas del tribunal}
\begin{enumerate}
\item \textbf{¿Hay desbordamientos sin ganancias de coordinar?} Sí: con LCP, B queda en $L$, pero los 2 eligen la misma política y el tipo de cambio no se mueve.
\item \textbf{¿Por qué cambia el signo de la monetaria entre M-F y C-P?} En M-F quita demanda al vecino; en C-P (PCP) le abarata lo que importa.
\item \textbf{¿Por qué la cooperativa no es un equilibrio?} Fuera de la función de reacción: cada uno tiene incentivo a desviarse.
\item \textbf{¿Qué aporta un foro?} Repite el juego ($\beta\uparrow$), informa, vigila, liga asuntos.
\item \textbf{¿Por qué no se coordinó tras 2010?} Incertidumbre sobre el modelo: unos jugaban a la disciplina, otros a la locomotora.
\item \textbf{G-20 frente a OCDE}: crisis y demanda, informal; normas y oferta, con tratado.
\item \textbf{¿Por qué coordinar puede ser contraproducente?} Si todos expanden, ninguna moneda se deprecia y desaparece un freno a sorprender con inflación (\autor{Rogoff}, 1985).
\item \textbf{¿La globalización pide más coordinación?} No necesariamente: la relación de intercambio y los mercados financieros ya reparten el riesgo (\autor{Bergin}, 2002).
\item \textbf{¿Qué es la guerra de divisas y cómo se paró?} Depreciar es dominante y, si todos lo hacen, nadie gana; G-7, 2013: objetivos internos con instrumentos internos.
\item \textbf{¿Quién ajusta, deudor o acreedor?} El deudor, cuando se le acaba la financiación; \autor{Keynes} quería cobrar también al acreedor; hoy, presión sin sanción (ESR, G-20, procedimiento de desequilibrios).
\item \textbf{¿Por qué el FSB y no el FEF?} El G-20 lo creó en Londres (4/2009) para coordinar a reguladores y supervisores tras la crisis; sucede al Foro de Estabilidad Financiera de 1999 \vertema{3.B.27}.
\item \textbf{¿Qué coordina la OCDE si no hace acuerdos de demanda?} Las políticas de oferta: normas (Código de capitales, antisoborno, multinacionales), \emph{Going for Growth}, el Consenso de créditos a la exportación y los 2 pilares.
\end{enumerate}

\begin{trampa}
\begin{itemize}
\item Fiscal con flexible = locomotora; monetaria con flexible = empobrecimiento (M-F). Con tipo fijo, la monetaria sube las 2 rentas.
\item LCP: \emph{sin} ganancias de coordinar; dólar: sí, unilaterales.
\item FSB: abril de 2009 (no 2013: en 2013 pasa a asociación suiza).
\item Plaza: G-5 (1985); Louvre: G-6 (1987). G-7 de 2013: declaración de ministros, no cumbre.
\item \autor{Rogoff} (1985), \autor{Oudiz}-\autor{Sachs} (1984); \autor{Koichi} Hamada, no \autor{Robert} Hamada.
\item Unión Africana en el G-20 desde 2023; España es invitado permanente, no miembro.
\end{itemize}
\end{trampa}

\end{columnas}
\end{document}
